\documentclass[11pt,a4paper]{article}

\usepackage[margin=1in]{geometry}
\usepackage{amsmath,amssymb}
\usepackage{booktabs}
\usepackage{graphicx}
\usepackage{hyperref}
\usepackage{natbib}
\usepackage{caption}
\usepackage{microtype}

\hypersetup{colorlinks=true, linkcolor=black, citecolor=black, urlcolor=blue}

\title{\textbf{Gridiron: A Market-Anchored Correlated Monte Carlo Model\\
for Fantasy Football Projection}\\[0.4em]
\large Design, empirical calibration, and a held-out season}

\author{S. Lifshitz}
\date{August 2026}

\begin{document}
\maketitle

\begin{abstract}
\noindent
We describe a simulation model for American football that produces joint
distributions over player fantasy scoring, rather than point estimates. Three
design choices distinguish it. First, scoring is sampled \emph{per game} --- one
margin and one total per fixture --- so that teammates and opponents share the
game state that generates their statistics, and correlations arise from the
structure rather than being imposed afterwards. Second, team totals are anchored
to closing betting-market lines, and every candidate adjustment is tested for
whether the market has already priced it; adjustments that it has are discarded
as double counting. Third, all effect sizes are estimated from data rather than
asserted, and effects that fail to separate from noise are reported as zero and
retained in the codebase as negative results.

We estimate environmental effects on 3{,}341 outdoor team-games and 12{,}787
field-goal attempts, and matchup effects out of sample on 56{,}554 player-weeks.
We find that wind is \emph{not} priced into closing totals ($-1.83$ points per
$10$ mph, significant) while temperature is; that altitude has no measurable
effect on field-goal accuracy once attempt distance is controlled
($-0.01$ yards per $1{,}000$ ft, $\mathrm{se}=0.028$); and that team-versus-team
history carries no information beyond the closing spread ($r=0.026$), whereas a
defence's history against a \emph{position} does ($\beta$ between $0.16$ and
$0.27$, all significant).

On a strictly held-out 2025 season --- 272 games and 5{,}237 player-weeks, with
all results and statistics from that season hidden from every estimator --- the
model matches the closing line on game outcomes (margin MAE $9.72$ against the
market's $9.72$) without beating it (50.7\% against the spread), and is well
calibrated on win probability (Brier $0.212$ against a $0.249$ base rate,
calibration error $0.023$). Player-level rank correlation with realised scoring
is $0.48$ overall and $0.67$ for running backs. We also document two defects:
projection intervals are too narrow, containing $69.3\%$ of outcomes at a nominal
$80\%$, and a residual under-projection of $-0.82$ points per player-week persists
after adjusting for availability.
\end{abstract}

\section{Introduction}

Fantasy football scoring is a deterministic function of a stochastic process.
Given the play-by-play record of a game, a player's fantasy points follow
mechanically from a scoring rule; all of the uncertainty lives upstream, in
which plays happen and to whom. This suggests an obvious modelling strategy ---
simulate the upstream process and apply the scoring rule --- and yet the
projections that dominate public fantasy tools are point estimates, typically a
single expected number per player per week.

A point estimate is the wrong object for almost every decision a fantasy manager
actually makes. Choosing between two flex options is a question about which
distribution is preferable given the rest of a lineup, not which mean is larger.
Deciding whether to trade a consistent producer for a volatile one is explicitly
a question about higher moments. Estimating a team's championship odds requires
the joint distribution across a roster, because a manager who stacks a
quarterback with his own receivers has a very different variance profile from
one who does not, even at identical expected points.

This paper describes \emph{Gridiron}, a system built around that observation.
Its output is a set of simulated seasons and weeks --- by default $20{,}000$ of
them --- from which any functional of the joint distribution can be read.

Three commitments shape the design, and we state them at the outset because they
are what the paper is really about.

\paragraph{Correlation should be structural, not bolted on.}
Two players in the same game are not independent draws. When a game turns into a
shootout, both quarterbacks benefit; when a team falls behind, its running back
loses carries and its receivers gain targets. Models that generate each player
independently and then apply a correlation matrix are approximating a structure
they could have simulated directly. We sample one margin and one total per
fixture, allocate points to teams, allocate team production to players, and let
the correlations fall out.

\paragraph{The market is a strong prior, and double counting is the failure mode.}
Closing betting lines aggregate more information than any freely available model
will reconstruct. Anchoring to them is straightforward. The subtlety is what
happens next: if the model then applies its own adjustment for a factor the
market has \emph{already} priced, it counts that factor twice and degrades a
good prior. We therefore test every candidate adjustment against the residual
from the closing line, not against raw outcomes, and apply only what survives.
Section~\ref{sec:environment} shows this distinction changing a design decision.

\paragraph{Measured, including when the measurement is zero.}
Every coefficient reported here was estimated from data. Where an effect that
practitioners believe in fails to separate from noise, we report zero and retain
the negative result in the source, so that the absence of a knob is legible as a
finding rather than an oversight. Two such results appear below: altitude on
kicking, and head-to-head history.

\section{Data}

All football data comes from the \texttt{nflverse} project, which publishes
cleaned play-by-play, weekly player statistics, team statistics, rosters, depth
charts, injury reports and schedules with closing betting lines. Consensus draft
rankings come from a public mirror of FantasyPros expert consensus. Venue
metadata --- coordinates, field elevation, roof type and surface for all 32
current stadiums plus international sites --- was compiled for this work, as
\texttt{nflverse} carries roof and surface per game but not location or
elevation. Forecast weather is retrieved from Open-Meteo, which requires no
credentials.

Estimation samples are stated with each result. The two largest are 3{,}341
outdoor team-games carrying both temperature and wind (1999--2025) and 12{,}787
field-goal attempts with kick distance (2014--2025).

\section{Model}

\subsection{Game state}

For each fixture $g$ we sample a margin and a total,
\begin{align}
M_g &\sim \mathcal{N}\!\left(s_g,\ \sigma_M^2\right), \\
T_g &\sim \mathcal{N}\!\left(t_g,\ \sigma_T^2\right),
\end{align}
where $s_g$ and $t_g$ are the closing spread and total, and the dispersions
$\sigma_M = 12.73$ and $\sigma_T = 13.19$ points are the residual standard
deviations of realised outcomes around closing lines, estimated on the same
historical sample. Team points follow by construction,
\begin{equation}
P^{\text{home}}_g = \tfrac{1}{2}(T_g + M_g), \qquad
P^{\text{away}}_g = \tfrac{1}{2}(T_g - M_g),
\end{equation}
clipped to $[0, 80]$.

Sampling at the level of the fixture rather than the team is the single most
consequential structural choice in the model. Drawing two independent team
scores permits both sides of a game to be simulated as winning, which is not a
state of the world; more importantly, it destroys precisely the correlation that
makes lineup construction interesting.

\subsection{Usage allocation}

Team production is allocated to players by Dirichlet draws over usage shares,
\begin{equation}
(\pi_{1}, \dots, \pi_{K}) \sim \mathrm{Dir}(\alpha \bar{\pi}_{1}, \dots, \alpha \bar{\pi}_{K}),
\end{equation}
where $\bar{\pi}_k$ are projected mean shares and the concentration $\alpha$ is
fitted to observed week-to-week share variance. Separate allocations govern
targets, carries and quarterback dropbacks.

Mean shares are estimated by empirical-Bayes shrinkage of each player's history
toward a positional and depth-chart prior, with explicit prior sample sizes so
that a player with three games of history is pulled hard toward the prior and
one with three seasons is not. Rates such as catch rate, yards per carry and
interception rate receive the same treatment.

\subsection{Season-level latent factors}

Independent weekly sampling produces a leaderboard that is far too flat: it
generates the average season repeatedly rather than the spread of seasons that
actually occurs. We therefore draw, once per simulated season per player, latent
multipliers on role and efficiency, together with a team-strength surprise
($\sigma \approx 2.6$ points per game) and a pass/rush touchdown split. These
persist across all weeks of that simulated season.

A subtlety arises here that is easy to get wrong. A log-normal multiplier on
Dirichlet concentration parameters inflates expected shares by Jensen's
inequality, drifting the entire board upward by roughly 10\%. We normalise the
shock by its share-weighted team mean so that expected shares are preserved.

\subsection{Availability}

Absences are persistent, not independent across weeks. We model availability as
a two-state Markov chain per player, with transition probabilities fitted from
injury-report history, so that a simulated injury produces a run of missed games
rather than isolated zeros.

\subsection{Market blending}

Bottom-up projections are blended with consensus expert rankings by isotonic
regression of model points-per-game on market rank, fitted \emph{within each
position}. Fitting across positions is incorrect: an overall draft board encodes
positional scarcity, so a quarterback taken 107th is not the 107th-best scorer,
and only the within-position ordering carries points information.

\section{Environmental effects}
\label{sec:environment}

\subsection{Identification}

Marginal comparisons of outcomes across weather conditions confound weather with
the teams that play in it: cold-weather franchises differ systematically from
dome franchises. We therefore estimate within-team fixed effects, demeaning both
outcomes and conditions by team-season, so that the identifying variation is the
same team playing in better or worse conditions than it usually does.

The second and more important choice is the dependent variable. To decide
whether an adjustment should be applied \emph{on top of} a market-anchored
total, the relevant quantity is not the outcome but the residual from the
closing line. An effect visible in raw scoring but absent from the residual is
one the market has already priced.

\subsection{Results}

Table~\ref{tab:weather} reports both. Per $10$ mph of wind, passing yards fall
by $16.3$, completion rate by $1.5$ percentage points and yards per attempt by
$0.32$; rushing share rises by $0.017$. Critically, the residual from the
closing total falls by $1.83$ points and is significant: \emph{the market does
not fully price wind}. Per $10^\circ$F of cold, the compositional effects are
present and significant but the residual is $-0.07$ points and is not.

\begin{table}[t]
\centering
\caption{Environmental effects, within-team fixed effects, 3{,}341 outdoor
team-games. ``vs.\ line'' is the residual of realised total from the closing
total --- the quantity that determines whether an adjustment is additive to a
market-anchored model or duplicative of it.}
\label{tab:weather}
\begin{tabular}{lrrl}
\toprule
& per $+10$ mph wind & per $-10^\circ$F & \\
\midrule
Rushing share        & $+0.0167$ & $+0.0071$ & both significant \\
Completion rate      & $-0.0152$ & $-0.0066$ & both significant \\
Yards per attempt    & $-0.324$  & $-0.082$  & both significant \\
Passing yards        & $-16.30$  & $-3.57$   & both significant \\
Rushing yards        & $+1.71$   & $+1.84$   & wind n.s. \\
\midrule
\textbf{Total vs.\ line} & $\mathbf{-1.83}$ & $-0.07$ & \textbf{wind significant}; cold n.s. \\
\bottomrule
\end{tabular}
\end{table}

The model therefore applies the wind effect to the game total and the
temperature effect to composition only. Applying temperature to the total as
well would count it twice, since the closing line already reflects it.

\subsection{Kicking, and a negative result on altitude}

Raw field-goal percentage appears to fall with wind and then recover above
$15$ mph, which is an artefact of attempt selection: coaches stop attempting
long kicks in difficult conditions, so the surviving attempts are easier. We
therefore fit make probability conditional on distance,
\begin{equation}
\Pr(\text{made}) = \sigma\!\left(\beta_0 + \beta_d d + \beta_w w + \beta_t \tau + \beta_e e\right),
\end{equation}
on 12{,}787 attempts. Results appear in Table~\ref{tab:kick}.

\begin{table}[t]
\centering
\caption{Field-goal accuracy conditional on distance, 12{,}787 attempts
(8{,}372 outdoor with weather). Coefficients are on the logit scale; the final
column converts to equivalent yards of kick distance.}
\label{tab:kick}
\begin{tabular}{lrrr}
\toprule
& coefficient & std.\ error & equivalent yards \\
\midrule
Distance (per yard)      & $-0.1116$ & $0.0038$ & --- \\
Wind (per $10$ mph)      & $-0.1817$ & $0.0626$ & $-1.63$ \\
Temperature (per $-10^\circ$F) & $-0.0186$ & $0.0188$ & n.s. \\
Elevation (per $1{,}000$ ft)   & $-0.0013$ & $0.0282$ & $-0.01$, n.s. \\
\bottomrule
\end{tabular}
\end{table}

Wind costs $1.63$ yards of effective range per $10$ mph. Altitude does not
separate from zero. Denver's attempts average $39.2$ yards against $38.7$
elsewhere, and conditional accuracy is flat in elevation. The underlying physics
is not in doubt --- air density at $5{,}280$ ft is roughly $82\%$ of sea level
--- but the effect on professional kickers is below the noise floor of this
sample. The model applies no altitude adjustment, and the estimate is retained
in the source so that its absence reads as a finding.

\section{Matchup effects}

Two quantities are commonly called ``matchup history'' and they behave
differently.

\paragraph{Team-versus-team history predicts nothing.}
On 5{,}798 games with at least two prior meetings, regressing the residual from
the closing spread on the mean margin of the last six meetings gives
$\beta = 0.043$ ($\mathrm{se} = 0.022$), $r = 0.026$, $R^2 = 0.0007$; using only
the most recent meeting gives $r = 0.003$. Neither is significant. Whatever it
means for one team to ``own'' another, it does not survive a point spread. The
system computes and displays this history, and weights it at exactly zero.

\paragraph{Defence-versus-position history does predict.}
Evaluated strictly out of sample --- prior weeks only, predicting the following
week --- a defence's relative fantasy points allowed to a position carries the
coefficients in Table~\ref{tab:matchup}, all significant.

\begin{table}[t]
\centering
\caption{Out-of-sample predictive weight of a defence's prior relative fantasy
points allowed, by position. The fitted $\beta$ \emph{is} the shrinkage applied
in production.}
\label{tab:matchup}
\begin{tabular}{lrrrr}
\toprule
Position & $n$ & $\beta$ & std.\ error & $r$ \\
\midrule
QB & 3{,}989 & $0.209$ & $0.040$ & $0.082$ \\
RB & 3{,}998 & $0.266$ & $0.039$ & $0.108$ \\
WR & 3{,}998 & $0.176$ & $0.045$ & $0.062$ \\
TE & 3{,}979 & $0.163$ & $0.041$ & $0.063$ \\
\bottomrule
\end{tabular}
\end{table}

These coefficients are not chosen; they \emph{are} the correct shrinkage. A
defence that has allowed $67\%$ more than average to tight ends should be
expected to allow about $11\%$ more next week, not $67\%$: regression to the
mean is most of the story, and $\beta$ measures how much of it. The adjustment
is applied weekly only. Over a full season a team faces seventeen defences and
the effect averages out, so incorporating it into a season-long draft ranking
would add noise in the costume of information.

\section{Held-out evaluation}

\subsection{Protocol}

We hold out the 2025 season entirely. A cutoff mechanism caps every historical
load at 2024, including explicit multi-season requests, so that usage profiles,
efficiency priors, team-strength estimates and all fitted constants are computed
without any 2025 observation. Schedules for 2025 remain visible --- fixtures and
closing lines are pre-game information and are the anchor the model is built on
--- but their result columns are redacted at load time, so no estimator can
observe an outcome.

We regard this redaction as the essential detail. An earlier version of the
protocol capped only the default ``load all history'' path, while
\texttt{team\_context} requested seasons explicitly and thereby saw the entire
holdout season. That version reported a margin MAE of $6.97$ against a market
benchmark of $12.03$ and an implied $84\%$ against the spread --- results that
are impossible, and whose implausibility is what exposed the leak. We report the
episode because a leaked backtest is worse than none, and because its signature
was not a warning but an implausibly good number.

Weekly matchup factors are permitted to use prior weeks within the holdout
season, as they would be in live use.

\subsection{Determinism as a correctness check}

Establishing the holdout also surfaced a class of defect we had not been looking
for. Running the backtest twice on identical inputs produced different metrics
--- margin MAE moving in the third decimal, and the against-the-spread rate
moving by six percentage points. Since the simulation is explicitly seeded, this
should have been impossible, and chasing it found three genuine bugs upstream of
the simulator.

First, the market board assigned integer ranks after sorting on average draft
position alone. Players tied on ADP therefore took ranks in whatever order the
sort happened to emit. Because the market blend is a global isotonic fit over
those ranks, a permutation of three tied players moved projected points for most
of the board.

Second, kickers were modelled by a dedicated routine \emph{and} carried in the
general roster pool, and the two were concatenated. Twenty-nine players appeared
on the board twice, where they double-counted in every team-level aggregation.

Third, the routine that reduces each team to one kicker did so with an
order-dependent deduplication. Beyond flipping between runs, it selected the
wrong player: Cincinnati resolved to a kicker with no recent attempts over the
incumbent starter. Selection now ranks by attempts in the previous two seasons
with an identifier tiebreak.

None of these would have been caught by the test suite, which asserts properties
of a single run. All three were caught by the weaker requirement that two
identical runs agree. We note it because determinism is usually filed under
engineering hygiene, and here it functioned as a correctness test that the
statistical checks did not replicate. After the fixes the backtest reproduces
bit-for-bit, and all figures below are from a reproducible run.

\subsection{Game-level results}

Table~\ref{tab:game} reports performance over 272 regular-season games.

\begin{table}[t]
\centering
\caption{Held-out 2025 season, game level, 272 games at $6{,}000$ simulations
each. Market benchmarks use the closing line.}
\label{tab:game}
\begin{tabular}{lrr}
\toprule
& Model & Closing line \\
\midrule
Margin MAE            & $9.71$  & $9.72$ \\
Margin RMSE           & $12.27$ & --- \\
Total MAE             & $10.39$ & $10.39$ \\
Total RMSE            & $13.18$ & --- \\
\midrule
Win accuracy          & $0.654$ & --- \\
Brier score           & $0.212$ & base rate $0.249$ \\
Log loss              & $0.608$ & --- \\
Calibration error     & $0.023$ & --- \\
Beat the closing line & $0.559$ & $z = 1.94$, n.s. \\
\bottomrule
\end{tabular}
\end{table}

The model reproduces the closing line almost exactly on both margin and total,
and does not demonstrably beat it. The $55.9\%$ against the spread is the one
figure here that might look like an edge, and it is not one: over 272 games the
standard error on a coin flip is $3.0$ percentage points, so this is $1.94$
standard errors from chance --- the kind of result that arises about one season
in twenty by luck alone. A single season cannot separate a $56\%$ bettor from a
$50\%$ one, and we do not claim to be the former. This is what a model should
look like when it takes the closing line as its prior and adjusts only for
effects it can demonstrate are unpriced. Win probabilities are well calibrated, with a calibration error of
$0.023$ and close agreement between predicted and realised frequencies in every
bin below $0.8$; the top bin ($n = 23$) realises $0.957$ against a predicted
$0.853$, which at that sample size is two games.

We emphasise this null result. The system is not a betting edge on game
outcomes, and reporting it as one would be false. Its value lies downstream, in
allocating a well-calibrated game distribution to individual players.

\subsection{Player-level results}

Table~\ref{tab:player} reports 5{,}237 player-weeks. Because a projection
averages over simulations in which a player was injured or benched and scored
zero, while an actual is observed only when he played, we report error both
unconditionally and divided by the simulated active share.

\begin{table}[t]
\centering
\caption{Held-out 2025 season, player level. ``$\mid$played'' divides the
projection by simulated availability, which is the like-for-like comparison
against outcomes that are observed only when the player appeared. The final
column is the share of outcomes inside the nominal $80\%$ interval.}
\label{tab:player}
\begin{tabular}{lrrrrrrr}
\toprule
& $n$ & MAE & MAE$\mid$played & bias & bias$\mid$played & Spearman $\rho$ & in 80\% band \\
\midrule
All & 5{,}237 & $4.97$ & $4.96$ & $-1.65$ & $-0.82$ & $0.482$ & $69.3\%$ \\
QB  & 564     & $7.12$ & $6.90$ & $-3.46$ & $-1.83$ & $0.445$ & $60.6\%$ \\
RB  & 1{,}317 & $4.63$ & $4.55$ & $-2.88$ & $-2.11$ & $0.666$ & $75.3\%$ \\
WR  & 2{,}047 & $4.39$ & $4.47$ & $-1.27$ & $-0.33$ & $0.638$ & $83.4\%$ \\
TE  & 929     & $4.22$ & $4.07$ & $-3.54$ & $-3.19$ & $0.541$ & $62.0\%$ \\
\bottomrule
\end{tabular}
\end{table}

Rank correlation with realised scoring is $0.67$ for running backs and $0.64$
for receivers, the two positions where usage is most persistent, and weakest for
quarterbacks, where a small population and high touchdown variance dominate.

\section{Limitations}

Two defects are visible in Table~\ref{tab:player} and we state them plainly.

\paragraph{Intervals are too narrow.}
The nominal $80\%$ interval contains $69.3\%$ of outcomes, and only $60.6\%$ for
quarterbacks and $62.0\%$ for tight ends. The model understates tail risk, which
means published boom and bust rates are conservative and the distribution is
sharper than the world. For a system whose entire premise is that the
distribution matters more than the mean, this is the most consequential
weakness, and closing it should take priority over any improvement to central
tendency.

\paragraph{Residual under-projection.}
After adjusting for availability, projections remain low by $0.82$ points per
player-week overall and by $3.19$ for tight ends. Part of this is selection: the
evaluation sample contains players who accumulated statistics, and preseason
projections spread opportunity across a depth chart while realised opportunity
concentrates. Part is likely genuine miscalibration of touchdown rates. We have
not separated the two.

Further limitations: the model has no representation of benching, which is why
replacement-level quarterbacks project too favourably; weather beyond
temperature, wind and forecast precipitation is ignored, and precipitation
effects are a stated prior rather than a fitted coefficient because
\texttt{nflverse} does not carry precipitation; kicker and defence models are
deliberately simple; and the evaluation covers a single season, so all
differences reported here are subject to one-season sampling variation.

\section{Conclusion}

The contribution of this work is less a modelling technique than a discipline.
Simulating games rather than players yields correlations for free. Anchoring to
market lines yields a strong prior for free. What is \emph{not} free, and what
we argue deserves more attention than it usually receives, is the question of
which adjustments the prior has already absorbed. Wind survives that test;
temperature does not. Positional matchup history survives it; team-versus-team
history does not. Altitude, despite sound physics and firm folk belief, does not
survive at any sample size available to us.

On a genuinely held-out season the resulting system matches the closing line
without beating it, produces well-calibrated win probabilities, and ranks player
scoring with a correlation of $0.48$ overall and $0.67$ for running backs. It
also produces intervals that are too narrow and projections that are slightly
too low. We report the negative results and the defects alongside the positive
ones, on the view that a projection system whose failures are documented is more
useful than one whose successes are.

\section*{Reproducibility}

The system, the fitting scripts and the backtest harness are available in the
project repository. The holdout is reproduced with
\texttt{python scripts/backtest.py --season 2025 --sims 6000}, which sets the
training cutoff, rebuilds every estimator from data through 2024, predicts all
272 games and writes the metrics reported here. All figures in
Tables~\ref{tab:weather}--\ref{tab:player} are produced by that pipeline.

\end{document}
